
\SGARefTarget{sga3:iv:chapter:iv}\chapter*{Expos\'e IV\\Topologies and Sheaves}
\addcontentsline{toc}{chapter}{Expos\'e IV. Topologies and Sheaves}
\markboth{EXPOS\'E IV. TOPOLOGIES AND SHEAVES}%
  {EXPOS\'E IV. TOPOLOGIES AND SHEAVES}

\begingroup
\renewcommand{\thefootnote}{\fnsymbol{footnote}}
\begin{center}
  \textit{by M. Demazure}\footnote{This text develops the substance of two
  lectures given by A.~Grothendieck, supplementing them at several important
  points that had been passed over in silence or only briefly touched upon.}
\end{center}
\endgroup


This expos\'e is intended to acquaint the reader with the essentials of the
language of topologies and sheaves (without cohomology), a language
particularly convenient for questions concerning the formation of quotients,
among others.

The first three sections develop the language of forming quotients. The
fourth, which is the central part, presents the theory of sheaves, oriented
mainly toward its application to questions about quotients; the fifth applies
this theory to the formation of quotients in groups and to principal
homogeneous bundles. The final section concerns the category of schemes more
specifically and defines several useful topologies on that category.

The reader will profitably consult \SGARefLink{sga3:iv:bib:as}{[AS]}, \SGARefLink{sga3:iv:bib:ma}{[MA]}, \SGARefLink{sga3:iv:bib:d}{[D]}, and \SGARefLink{sga3:iv:bib:sga:4}{SGA~4}; in particular,
\SGARefLink{sga3:iv:bib:d}{[D]} for the applications of topologies to descent theory, and \SGARefLink{sga3:iv:bib:sga:4}{SGA~4} for
questions concerning universes (which are treated especially poorly in this
expos\'e).

\SGARefTarget{sga3:iv:section:1}\section*{1. Universal Effective Epimorphisms}
\addcontentsline{toc}{section}{1. Universal Effective Epimorphisms}

Throughout the remainder of this expos\'e, a category \(\mathscr C\) is fixed.

\medskip
\noindent\SGARefTarget{sga3:iv:definition:1:1}\textbf{Definition 1.1.}
\textit{A morphism \(u:T\to S\) is called an \emph{epimorphism} if, for every
object \(X\), the corresponding map}
\[
  X(S)=\operatorname{Hom}(S,X)
  \longrightarrow
  X(T)=\operatorname{Hom}(T,X)
\]
\textit{is injective.}
\textit{It is called a \emph{universal epimorphism} if, for every morphism
\(S'\to S\), the fiber product \(T'=T\times_S S'\) exists and
\(u':T'\to S'\) is an epimorphism.}

\medskip
\noindent\SGARefTarget{sga3:iv:definition:1:2}\textbf{Definition 1.2.}
\textit{A diagram of maps of sets}
\[
  A\xrightarrow{u}B
  \overset{v_1}{\underset{v_2}{\rightrightarrows}}C
\]
\textit{is said to be \emph{exact} if \(u\) is injective and its image consists
of those elements \(b\in B\) for which \(v_1(b)=v_2(b)\). A diagram of the
same form in \(\mathscr C\) is said to be \emph{exact} if, for every object
\(X\) of \(\mathscr C\), the corresponding diagram of sets}
\[
  A(X)\longrightarrow B(X)\rightrightarrows C(X)
\]
\textit{is exact; one then also says that \(u\) makes \(A\) a \emph{kernel}
of the pair of arrows \((v_1,v_2)\). Dually, a diagram in \(\mathscr C\)}
\[
  C\overset{v_1}{\underset{v_2}{\rightrightarrows}}B
  \xrightarrow{u}A
\]
\textit{is said to be exact if it is exact as a diagram in the opposite
category \(\mathscr C^{\circ}\); that is, if for every object \(X\) of
\(\mathscr C\), the corresponding diagram of sets}
\[
  X(A)\longrightarrow X(B)\rightrightarrows X(C)
\]
\textit{is exact.}
\textit{One also says that \(u\) makes \(A\) a \emph{cokernel} of the pair of
arrows \((v_1,v_2)\).}

\medskip
\noindent\SGARefTarget{sga3:iv:definition:1:3}\textbf{Definition 1.3.}
\textit{A morphism \(u:T\to S\) is called an \emph{effective epimorphism} if
the fiber square \(T\times_S T\) exists and the diagram}
\[
  T\times_S T
  \overset{\operatorname{pr}_1}{\underset{\operatorname{pr}_2}{\rightrightarrows}}
  T\xrightarrow{u}S
\]
\textit{is exact; that is, if \(u\) makes \(S\) a cokernel of
\((\operatorname{pr}_1,\operatorname{pr}_2)\). It is called a
\emph{universal effective epimorphism} if, for every morphism \(S'\to S\),
the fiber product \(T'=T\times_S S'\) exists and the morphism
\(u':T'\to S'\) is an effective epimorphism.}

There are evidently implications
\[
\begin{array}{ccc}
 \text{universal effective epimorphism}
   &\Longrightarrow& \text{effective epimorphism}\\
 \Big\Downarrow && \Big\Downarrow\\
 \text{universal epimorphism}
   &\Longrightarrow& \text{epimorphism},
\end{array}
\]
but in general no other implication holds.\footnote{Editor's note: for
example, if \(\mathscr C=\mathbf{Sch}\) is the category of schemes, it is
easy to see that every universal epimorphism is surjective. Let
\[
  T=\coprod_{p\ \mathrm{prime}}\operatorname{Spec}(\mathbf F_p),
  \qquad S=\operatorname{Spec}(\mathbf Z).
\]
Then the morphism \(u:T\to S\) is an epimorphism that is not universal. On
the other hand, \(T\times_S T\) identifies with \(T\), so the two
projections \(T\times_S T\to T\) coincide. Since
\(\operatorname{id}_T\) does not descend to a morphism \(S\to T\), this
shows that \(u\) is not an effective epimorphism.}

\medskip
\noindent\SGARefTarget{sga3:iv:definition:1:4:0}\textbf{Definition 1.4.0.}
We ``recall'' that a morphism \(u:T\to S\) is said to be
\emph{squarable} if, for every morphism \(S'\to S\), the fiber product
\(T\times_S S'\) exists.

\medskip
\noindent\SGARefTarget{sga3:iv:lemma:1:4}\textbf{Lemma 1.4.}
\textit{Consider morphisms \(U\xrightarrow{v}T\xrightarrow{u}S\). Then:}
\begin{enumerate}
  \item[(a)]\SGARefTarget{sga3:iv:lemma:1:4:item:a:01} \textit{if \(u\) and \(v\) are epimorphisms, then \(uv\) is an
    epimorphism; and if \(uv\) is an epimorphism, then \(u\) is an
    epimorphism;}
  \item[(b)]\SGARefTarget{sga3:iv:lemma:1:4:item:b:01} \textit{if \(u\) and \(v\) are universal epimorphisms, then
    \(uv\) is a universal epimorphism; and if \(uv\) is a universal
    epimorphism and \(u\) is squarable, then \(u\) is a universal
    epimorphism.}
\end{enumerate}

Lemma~\SGARefLink{sga3:iv:lemma:1:4}{1.4} follows immediately from the definitions. We deduce:

\medskip
\noindent\SGARefTarget{sga3:iv:corollary:1:5}\textbf{Corollary 1.5.}
\textit{Let \(u:X\to Y\) and \(u':X'\to Y'\) be universal epimorphisms,
and suppose that \(Y\times Y'\) exists. Then \(X\times X'\) exists and}
\[
  u\times u':X\times X'\longrightarrow Y\times Y'
\]
\textit{is a universal epimorphism.}

We also record the following.

\medskip
\noindent\SGARefTarget{sga3:iv:definition:1:6:0}\textbf{Definition 1.6.0.}
An object \(S\) of \(\mathscr C\) is called \emph{squarable} if its product
with every object of \(\mathscr C\) exists. (If \(\mathscr C\) has a final
object \(e\), this is equivalent to saying that the morphism \(S\to e\) is
squarable; cf. \SGARefLink{sga3:iv:definition:1:4:0}{1.4.0}.)

\medskip
\noindent\SGARefTarget{sga3:iv:lemma:1:6}\textbf{Lemma 1.6.}
\textit{Let \(u:X\to Y\) be a morphism in \(\mathscr C_{/S}\). For it to be
an epimorphism (respectively, a universal epimorphism, an effective
epimorphism, or a universal effective epimorphism), it is sufficient that the
corresponding morphism in \(\mathscr C\) have the same property; it is also
necessary if \(S\) is assumed to be a squarable object of \(\mathscr C\).}

The proof is immediate and is left to the reader. The hypothesis that \(S\)
is squarable is used to interpret the \(\mathscr C\)-morphisms from an object
\(Y\) of \(\mathscr C_{/S}\) to an object \(Z\) of \(\mathscr C\) as the
\(\mathscr C_{/S}\)-morphisms from \(Y\) to \(Z\times S\).
